% !TEX root = ../nc-cubic-non-as.tex
\section{Clifford order と内在的な三次曲面構造}
\label{sec:construction}

基礎体を $k=\mathbb C$ とし，加群は右加群とする．
Hirzebruch 曲面 $Z=\mathbb F_1$ の負切断を $C$，
射影 $\rho:Z\to\mathbb P^1$ のファイバー類を $f$ と書く．従って
\[
 C^2=-1,\qquad C\cdot f=1,\qquad f^2=0,
 \qquad K_Z=-2C-3f.
\]
次の線束と階数3のベクトル束を固定する．
\begin{equation}
 L=\mathcal O_Z(C+2f),\qquad
 \mathcal V=\mathcal O_Z^{\oplus2}\oplus\mathcal O_Z(f).
 \label{eq:choice-LV}
\end{equation}
$\mathbb F_1$ 上の ample cone の記述から $L$ は ample であり，$L^2=3$ である．また
\begin{equation}
 \det\mathcal V=\omega_Z\otimes L^{\otimes2}
 \label{eq:determinant-V}
\end{equation}
が成り立つ．

\subsection{二次形式と判別曲線}

$Z$ の Cox 座標を $s,t,c,d$ とし，その次数を
\[
 \deg s=\deg t=f,\qquad \deg c=C,\qquad \deg d=C+f
\]
とする．不安定集合を除いた座標空間は
$\{(s,t)\ne(0,0),\ (c,d)\ne(0,0)\}$ である．
$s,t$ のいずれかと $c,d$ のいずれかをそれぞれ1に正規化すると，
$Z$ を被覆する4つの標準アフィンチャートを得る．

$L$ に値を持つ対称形式 $q:\operatorname{Sym}^2\mathcal V\to L$ を
\begin{equation}
 q=
 \begin{pmatrix}
 cs^2+dt&ct^2&d\\
 ct^2&ct^2+ds&0\\
 d&0&c
 \end{pmatrix}
 \label{eq:explicit-q}
\end{equation}
で定める．第3行・列の束の次数は $f$ なので，
例えば $q_{13}$ の次数は $L-f=C+f$，$q_{33}$ の次数は $L-2f=C$ であり，
この行列は指定した束の間の射を定める．行列式は
\begin{equation}
 \begin{split}
 \Delta={}&c^3s^2t^2-c^3t^4+c^2ds^3+c^2dt^3\\
          &+cd^2st-cd^2t^2-d^3s.
 \end{split}
 \label{eq:explicit-discriminant}
\end{equation}
判別曲線を $D=(\Delta=0)$ と書く．その因子類は
$3L-2c_1(\mathcal V)=3C+4f$ である．

\begin{lemma}
\label{lem:smooth-discriminant}
$D$ は滑らかで連結な種数3の曲線である．
$q$ の階数は全ての点で2以上であり，付随する conic bundle
\[
 p:X=\{q=0\}\subset\mathbb P_Z(\mathcal V)\longrightarrow Z
\]
は flat で，全空間 $X$ は滑らかである．
ここで $\mathbb P_Z(\mathcal V)$ は $\mathcal V$ の直線をパラメータ化する．
\end{lemma}

\begin{proof}
各アフィンチャートで式\eqref{eq:explicit-discriminant}を正規化した多項式を
$F(u,v)$ とすると，特異点イデアル
$(F,\partial_uF,\partial_vF)\subset\mathbb Q[u,v]$ の辞書式
Gr\"obner 基底は次の通りである．
\[
\begin{array}{c|cccc}
\text{正規化}&s=c=1&s=d=1&t=c=1&t=d=1\\\hline
\text{Gr\"obner 基底}&[1]&[1]&[1]&[1]
\end{array}
\]
この計算の入力と検証方法は第\ref{sec:reproducibility}節に記す．
従って $D$ は $\mathbb C$ 上で滑らかである．
$K_Z+D=C+f$ と $H^1(Z,\mathcal O_Z(C+f))=0$ より，
Serre 双対性は $H^1(Z,\mathcal O_Z(-D))=0$ を与える．
因子の完全列から $H^0(D,\mathcal O_D)=k$ であり，$D$ は連結である．
また
\[
 2g(D)-2=D(D+K_Z)=(3C+4f)(C+f)=4
\]
だから $g(D)=3$ である．

行列の階数が1以下の点では全ての $2\times2$ 小行列式が消えるため，
行列式の微分も消える．これは $D$ の滑らかさに反する．
従って各ファイバーは非零な二次式で定義される conic であり，
Hilbert 多項式の一定性から $p$ は flat である．
$D$ 外のファイバーは滑らかである．$D$ 上では階数が2であり，
退化 conic の特異点における基底方向の微分が
判別式の非零微分に対応するので，$X$ も滑らかである．
\end{proof}

\subsection{Clifford order と一般点の階数}

Clifford 関係を $xy+yx=2q(x,y)$ と規約し，
偶部分と奇部分をそれぞれ
\[
 \mathcal A=\operatorname{Cl}_0(\mathcal V,q,L),\qquad
 \mathcal P=\operatorname{Cl}_1(\mathcal V,q,L)
\]
と書く．反対称化による束としての分解は
\begin{equation}
 \begin{split}
 \mathcal A
  &=\mathcal O_Z\oplus
     (\wedge^2\mathcal V)\otimes L^{-1}
    =\mathcal O_Z\oplus\mathcal O_Z(-C-2f)
       \oplus\mathcal O_Z(-C-f)^{\oplus2},\\
 \mathcal P
  &=\mathcal V\oplus(\det\mathcal V)\otimes L^{-1}
    =\mathcal O_Z^{\oplus2}\oplus\mathcal O_Z(f)
       \oplus\mathcal O_Z(-C-f)
 \end{split}
 \label{eq:underlying-bundles}
\end{equation}
である．積は Clifford 積であり，$\mathcal P$ は
$\mathcal A$-bimodule となる．以下
\[
 \mathcal C=\operatorname{coh}\mathcal A,\qquad
 \mathcal O=\mathcal A
\]
と置く．$\mathcal C$ は noetherian な $k$-線形アーベル圏である．

\begin{proposition}
\label{prop:division-degree-three}
一般点の代数 $\mathcal A\otimes_{\mathcal O_Z}k(Z)$ は
quaternion division algebra である．従って構造対象 $\mathcal O$ は
torsion-free で Gabriel--Krull rank 1 を持つ．
\end{proposition}

\begin{proof}
$i:D\hookrightarrow Z$ とし，
$q:\mathcal V\to\mathcal V^\vee\otimes L$ の余核を $i_*\mathcal M$ と書く．
$D$ 上での corank は1なので $\mathcal M$ は線束である．
$\mathcal K=\ker(q|_D)$，$\mathcal W=\mathcal V|_D/\mathcal K$ と置くと，
$\mathcal W$ は $L|_D$ に値を持つ非退化な階数2の二次束である．
その判別線束
\[
 \eta=\det\mathcal W\otimes L^{-1}|_D
     =\mathcal M\otimes\omega_Z|_D
\]
は $\eta^{\otimes2}\simeq\mathcal O_D$ を満たす．
符号を含めて $-\det(q|_{\mathcal W})$ を用いると，
$\mathcal O_D\oplus\eta$ は退化 conic の二つの成分を区別する
非分岐二重被覆の代数となる．
この線束には完全列
\begin{equation}
 0\longrightarrow\mathcal V\otimes\omega_Z
 \xrightarrow{q}\mathcal V^\vee\otimes L\otimes\omega_Z
 \longrightarrow i_*\eta\longrightarrow0
 \label{eq:eta-resolution-degree-three}
\end{equation}
がある．左の二つの束はそれぞれ
\[
 \mathcal O_Z(-2C-3f)^{\oplus2}\oplus\mathcal O_Z(-2C-2f),
 \qquad
 \mathcal O_Z(-C-f)^{\oplus2}\oplus\mathcal O_Z(-C-2f).
\]
後者はファイバー上の次数が $-1$ なので全コホモロジーが消える．
前者の $H^1$ は Serre 双対性と
$H^1(Z,\mathcal O_Z)=H^1(Z,\mathcal O_Z(-f))=0$ により消える．
従って $H^0(D,\eta)=0$ であり，$\eta$ は非自明である．
特に二重被覆 $\widetilde D\to D$ は連結である．

$D$ の一般点の離散付値環で $q=\langle a,b,uc\rangle$ と対角化する．
ここで $a,b,c$ は単元，$u$ は一様化元である．
偶 Clifford 代数の生成元 $I=e_1e_2,J=e_1e_3$ は
\[
 I^2=-ab,\qquad J^2=-auc,\qquad IJ=-JI
\]
を満たす．その Brauer residue は $\sqrt{-ab}$ による二次拡大であり，
上の二重被覆の関数体拡大と一致する．従って residue は非自明で，
一般点の quaternion algebra は split でない．次数2の中心単純代数は
split でなければ division である．
中心上有限なorderについて，定義\ref{def:weak-cubic}の
$\mathcal C_0$は零次元に台を持つ対象からなり，
$\mathcal C_1$は台の次元が一以下の対象からなる．
実際，零次元coherent加群は有限長であり，曲線に台を持つcoherent加群を
零次元対象で割ると，有限個の曲線の一般点での有限長加群に帰着する．
従って $\mathcal C/\mathcal C_1$ は，
一般点の division algebra の有限次元右加群圏である．
その正則加群の長さは一であり，$\mathcal A$は中心上局所自由なので
低次元に台を持つ部分加群を持たない．これで結論を得る．
\end{proof}

\subsection{局所正則性と反標準自己同値}

閉点の完備局所環を $B=k[[u,v]]$ とする．$D$ 外では
$\widehat{\mathcal A}\simeq M_2(B)$ である．$D=(u)$ 上では
単元の平方根を取って Clifford の偶生成元を正規化すると
$I^2=1,J^2=u,IJ=-JI$ とでき，
\begin{equation}
 \widehat{\mathcal A}\simeq
 \begin{pmatrix}B&B\\uB&B\end{pmatrix}
 \label{eq:local-order-degree-three}
\end{equation}
を得る．これは $u=r^2$ と置いた
$k[[r,v]]\rtimes\mu_2$ と同型であり，$\mu_2$ は $r$ の符号を変える．
群の位数が可逆なので誘導加群の平均化による分裂を用いると，
この skew group algebra の大域次元は2である．
実際，$k[[r,v]]$ 上の長さ2の射影分解を誘導して上界を得，
閉点加群の射影次元から下界を得る．
忠実平坦な完備化から射影性が降下するので，元の局所 order でも
coherent module の射影次元は2以下である．

各点の近傍で $q$ は単元を表すので，可逆な Clifford の奇元を取ることができる．
従って Clifford 積は bimodule の同型
\begin{equation}
 \mathcal P\otimes_{\mathcal A}\mathcal P\simeq\mathcal A\otimes L
 \label{eq:odd-square}
\end{equation}
を与える．また最高外冪成分への完全 pairing
\[
 \mathcal A\otimes\mathcal P\longrightarrow
       \det\mathcal V\otimes L^{-1}
\]
は偶元と奇元に関して cyclic なので，
$\mathcal A^\vee\simeq
\mathcal P\otimes(\det\mathcal V)^{-1}\otimes L$ は bimodule の同型である．
式\eqref{eq:determinant-V}より
\[
 \omega_{\mathcal A}:=
 \mathcal Hom_{\mathcal O_Z}(\mathcal A,\omega_Z)
 \simeq\mathcal P\otimes L^{-1}\simeq\mathcal P^{-1}.
\]
局所射影分解と $Z$ 上の双対性から，Serre functor は
\begin{equation}
 S=(-\otimes_{\mathcal A}\mathcal P^{-1})[2],\qquad
 T=S^{-1}[2]=-\otimes_{\mathcal A}\mathcal P,
 \qquad T^2=-\otimes L
 \label{eq:Serre-degree-three}
\end{equation}
である．従って $T$ はアーベル圏の完全自己同値である．
Serre 双対性と標準$t$-構造により，$i>2$なら
$\operatorname{Ext}^i_{\mathcal C}(M,N)=0$である．一方
\[
 \operatorname{Ext}^2_{\mathcal C}(\mathcal A,\mathcal P^{-1})
 \simeq\operatorname{End}_{\mathcal C}(\mathcal P^{-1})^*\ne0
\]
だから，$\mathcal C$のhomological dimensionは正確に二である．
Hom と Ext の有限次元性は中心 $Z$ の射影性から従う．
さらに conic bundle の半直交分解\cite[Theorem~4.2]{KuznetsovQuadric}
\begin{equation}
 D^b(X)=\left\langle
 D^b(Z,\mathcal A),\ p^*D^b(Z)\otimes\mathcal O_{X/Z}(1)
 \right\rangle
 \label{eq:conic-sod-degree-three}
\end{equation}
により，$D^b(\mathcal C)$ の標準 dg enhancement は smooth proper である．

\begin{proposition}
\label{prop:ample-degree-three}
$(\mathcal O,T)$ は Artin--Zhang ample である．反標準環
\[
 R=\bigoplus_{n\ge0}\operatorname{Hom}_{\mathcal C}
                  (\mathcal O,T^n\mathcal O)
\]
について $\mathcal C\simeq\operatorname{qgr}R$ が成り立つ．
\end{proposition}

\begin{proof}
任意の coherent $\mathcal A$-module $M$ は中心上でも coherent である．
$m\gg0$ では $M\otimes L^m$ は大域生成され，随伴により
$\mathcal A^{\oplus N}\twoheadrightarrow M\otimes L^m$ を得る．
これは $(T^{-2m}\mathcal O)^{\oplus N}\twoheadrightarrow M$ を与える．
次に全射 $M\twoheadrightarrow N$ の核を $K$ とし，$n=2m+r$，$r=0,1$ と書く．
$\mathcal A$ の局所射影性と Serre 消滅により
\[
 \operatorname{Ext}^1_{\mathcal A}(\mathcal A,T^nK)
 =H^1\bigl(Z,(K\otimes_{\mathcal A}\mathcal P^r)\otimes L^m\bigr)=0
 \qquad(m\gg0).
\]
これが ample pair のもう一つの条件である．
再構成は Artin--Zhang の定理\cite[Theorem~4.5]{AZ1994}による．
\end{proof}

\subsection{全次数の Hilbert 関数とコホモロジー消滅}

\begin{proposition}
\label{prop:Hilbert-degree-three}
構造対象 $\mathcal O$ は exceptional であり，全ての整数 $n$ に対して
\[
 \operatorname{Ext}^1_{\mathcal C}(\mathcal O,T^n\mathcal O)=0
\]
が成り立つ．$n<0$ では
$\operatorname{Hom}_{\mathcal C}(\mathcal O,T^n\mathcal O)=0$，
$n\ge0$ では
\begin{equation}
 \dim_k R_n=1+\frac32n(n+1),\qquad
 \operatorname{Ext}^{i}_{\mathcal C}(\mathcal O,T^n\mathcal O)=0
 \quad(i>0).
 \label{eq:Hilbert-degree-three}
\end{equation}
特に
\[
 H_R(t)=\frac{1-t^3}{(1-t)^4}.
\]
\end{proposition}

\begin{proof}
射影 $\rho$ に沿った計算により，$a\ge0$ に対して
\[
 \rho_*\mathcal O_Z(aC+bf)
       =\bigoplus_{i=0}^a\mathcal O_{\mathbb P^1}(b-i),
 \qquad R^1\rho_*\mathcal O_Z(aC+bf)=0.
\]
従って $b\ge a$ なら
\[
 h^0\mathcal O_Z(aC+bf)
   =\sum_{i=0}^a(b-i+1)
   =(a+1)(b+1)-\frac{a(a+1)}2,
\]
かつ高次コホモロジーは消える．
$a=-1$ の線束はファイバー上で acyclic なので全コホモロジーが消える．

式\eqref{eq:odd-square}と\eqref{eq:underlying-bundles}から
\[
 R_{2m}=H^0(Z,\mathcal A\otimes L^m),\qquad
 R_{2m+1}=H^0(Z,\mathcal P\otimes L^m).
\]
上の線束の公式を各直和因子に適用すると，$m\ge0$ の全範囲で
\begin{equation}
 \dim R_{2m}=6m^2+3m+1,\qquad
 \dim R_{2m+1}=6m^2+9m+4
 \label{eq:Hilbert-parities}
\end{equation}
を得る．$m=0$ で現れる負の線束はファイバー次数 $-1$ であり，
この場合も公式は成立する．同時に高次コホモロジーの消滅を得る．
$m<0$ では偶数・奇数のいずれの直和因子もファイバー次数が負なので，
全ての $n<0$ に対して $\operatorname{Hom}(\mathcal O,T^n\mathcal O)=0$ となる．
Serre 双対性は
\[
 \operatorname{Ext}^1(\mathcal O,T^n\mathcal O)^*
 \simeq\operatorname{Ext}^1(\mathcal O,T^{-n-1}\mathcal O)
\]
を与えるため，負の $n$ に対する $\operatorname{Ext}^1$ の消滅も従う．
$n=0$ の場合が構造対象の exceptionality である．
\end{proof}

\subsection{Hochschild homology}

\begin{proposition}
\label{prop:HH-degree-three}
標準的な dg enhancement について，HKR の次数を $p-q=i$ と規約すると
\[
 \operatorname{HH}_{-1}(\mathcal C)=k^2,\qquad
 \operatorname{HH}_0(\mathcal C)=k^4,\qquad
 \operatorname{HH}_1(\mathcal C)=k^2,
 \qquad \operatorname{HH}_i(\mathcal C)=0\quad(|i|\ge2)
\]
である．特に $D^b(\mathcal C)$ は full exceptional collection を持たない．
\end{proposition}

\begin{proof}
滑らかな判別曲線を持つ conic bundle に対する公式
\cite[Theorem~8.12]{KuznetsovHH}を上の次数規約に直すと
\[
 \operatorname{HH}_i(\mathcal C)
 \simeq\operatorname{HH}_i(Z)
 \oplus H^{-i}(D,\eta)
 \oplus H^{1-i}(D,\omega_D\otimes\eta)
\]
となる．$Z=\mathbb F_1$ の Hodge cohomology は対角部分だけで，
$\operatorname{HH}_0(Z)=k^4$ である．
$\eta$ は種数3の曲線上の非自明な2-torsion 線束なので，
Riemann--Roch と Serre 双対性は
\[
 h^1(D,\eta)=h^0(D,\omega_D\otimes\eta)=2,
 \qquad h^0(D,\eta)=h^1(D,\omega_D\otimes\eta)=0
\]
を与える．これで主張を得る．
full exceptional collection が存在すれば Hochschild homology の加法性により
非零次数の部分が消えるため，最後の結論も従う．
\end{proof}
